\documentclass[12pt,a4paper]{article}

\usepackage[utf8]{inputenc}
\usepackage[T1]{fontenc}
\usepackage{amsmath}
\usepackage{amssymb}
\usepackage{bm}
\usepackage{booktabs}
\usepackage{geometry}
\usepackage{hyperref}
\usepackage{listings}
\usepackage{xcolor}

\geometry{margin=2.5cm}

\definecolor{codebackground}{rgb}{0.95,0.95,0.95}
\lstdefinestyle{shellstyle}{
    backgroundcolor=\color{codebackground},
    basicstyle=\ttfamily\small,
    breaklines=true,
    frame=single,
    keepspaces=true,
    showstringspaces=false,
    tabsize=2
}
\lstset{style=shellstyle}

\title{Complex MaskNet pentru Speech Enhancement\\
\large Definitii, arhitectura si pipeline de rulare}
\author{}
\date{}

\begin{document}

\maketitle
\tableofcontents
\newpage

\section{Scopul modelului}

Modelul Complex MaskNet este o extensie a modelului MaskNet clasic folosit pentru
reducerea zgomotului din semnale de vorbire. Diferenta principala este ca modelul
nu mai invata doar o masca reala aplicata pe magnitudinea spectrogramei, ci invata
o masca complexa aplicata direct pe spectrograma complexa a semnalului zgomotos.

In varianta clasica, semnalul curatat este reconstruit folosind magnitudinea
estimata de retea si faza semnalului zgomotos. In varianta Complex MaskNet, reteaua
poate modifica atat componenta reala, cat si componenta imaginara a spectrogramei.
Astfel, reconstructia finala include implicit si o corectie a fazei.

\section{Modelul semnalului}

Fie semnalul zgomotos:
\begin{equation}
    y[n] = s[n] + d[n],
\end{equation}
unde:
\begin{itemize}
    \item $s[n]$ este semnalul curat de vorbire;
    \item $d[n]$ este zgomotul adaugat;
    \item $y[n]$ este semnalul observat, folosit ca intrare in model.
\end{itemize}

Semnalele sunt transformate in domeniul timp-frecventa prin STFT:
\begin{equation}
    Y[k,t] = \mathrm{STFT}\{y[n]\},
    \qquad
    S[k,t] = \mathrm{STFT}\{s[n]\},
\end{equation}
unde $k$ reprezinta indicele de frecventa, iar $t$ reprezinta cadrul temporal.

Spectrograma complexa se scrie:
\begin{equation}
    Y[k,t] = Y_r[k,t] + jY_i[k,t],
    \qquad
    S[k,t] = S_r[k,t] + jS_i[k,t].
\end{equation}

\section{De la MaskNet clasic la Complex MaskNet}

\subsection{MaskNet clasic}

MaskNet clasic primeste ca intrare magnitudinea spectrogramei zgomotoase:
\begin{equation}
    |Y[k,t]|.
\end{equation}

Reteaua prezice o masca reala:
\begin{equation}
    M_{\mathrm{mag}}[k,t] \in [0,1].
\end{equation}

Magnitudinea estimata a semnalului curat este:
\begin{equation}
    |\hat{S}[k,t]| = M_{\mathrm{mag}}[k,t] \cdot |Y[k,t]|.
\end{equation}

Pentru reconstructie, modelul refoloseste faza semnalului zgomotos:
\begin{equation}
    \hat{S}[k,t] =
    |\hat{S}[k,t]| \cdot e^{j \angle Y[k,t]}.
\end{equation}

Aceasta abordare este simpla si stabila, dar limiteaza calitatea reconstructiei
deoarece faza nu este estimata explicit.

\subsection{Complex MaskNet}

Complex MaskNet primeste doua canale de intrare:
\begin{equation}
    X_{\mathrm{in}}[k,t] =
    \begin{bmatrix}
        Y_r[k,t] \\
        Y_i[k,t]
    \end{bmatrix}.
\end{equation}

Reteaua prezice o masca complexa cu doua canale:
\begin{equation}
    M[k,t] = M_r[k,t] + jM_i[k,t].
\end{equation}

Spectrograma complexa curatata este obtinuta prin inmultire complexa:
\begin{equation}
    \hat{S}[k,t] = M[k,t] \cdot Y[k,t].
\end{equation}

Separand componentele reala si imaginara:
\begin{align}
    \hat{S}_r[k,t] &= M_r[k,t]Y_r[k,t] - M_i[k,t]Y_i[k,t], \\
    \hat{S}_i[k,t] &= M_r[k,t]Y_i[k,t] + M_i[k,t]Y_r[k,t].
\end{align}

Semnalul audio final se obtine cu iSTFT:
\begin{equation}
    \hat{s}[n] = \mathrm{iSTFT}\{\hat{S}[k,t]\}.
\end{equation}

\section{Complex Ratio Mask}

Target-ul folosit la antrenare este \textit{Complex Ratio Mask} (CRM). Acesta este
definit ca raportul complex dintre spectrograma curata si spectrograma zgomotoasa:
\begin{equation}
    M^{\star}[k,t] = \frac{S[k,t]}{Y[k,t] + \epsilon},
\end{equation}
unde $\epsilon$ este o constanta mica folosita pentru stabilitate numerica.

In implementare, raportul complex se calculeaza prin inmultirea cu conjugata:
\begin{equation}
    M^{\star}[k,t] =
    \frac{S[k,t]\overline{Y[k,t]}}{|Y[k,t]|^2 + \epsilon}.
\end{equation}

Pentru:
\begin{equation}
    S = S_r + jS_i,
    \qquad
    Y = Y_r + jY_i,
\end{equation}
componentele mastii complexe sunt:
\begin{align}
    M_r^{\star} &=
    \frac{S_rY_r + S_iY_i}{Y_r^2 + Y_i^2 + \epsilon}, \\
    M_i^{\star} &=
    \frac{S_iY_r - S_rY_i}{Y_r^2 + Y_i^2 + \epsilon}.
\end{align}

Valorile CRM pot deveni mari in zone unde energia semnalului zgomotos este foarte
mica. Pentru stabilitate, target-ul este limitat intr-un interval simetric:
\begin{equation}
    M^{\star} \leftarrow
    \mathrm{clip}(M^{\star}, -c, c),
\end{equation}
unde in configuratia curenta:
\begin{equation}
    c = 5.0.
\end{equation}

\section{Reprezentarea tensorilor}

Pentru antrenare, spectrograma complexa este convertita intr-un tensor real cu
doua canale:
\begin{equation}
    \bm{X} \in \mathbb{R}^{B \times 2 \times K \times T},
\end{equation}
unde:
\begin{itemize}
    \item $B$ este dimensiunea batch-ului;
    \item $2$ reprezinta canalele real si imaginar;
    \item $K$ este numarul de bin-uri de frecventa;
    \item $T$ este numarul de cadre temporale.
\end{itemize}

In proiect, pentru $n_{\mathrm{fft}}=512$:
\begin{equation}
    K = \frac{512}{2} + 1 = 257.
\end{equation}

Intrarea modelului este:
\begin{equation}
    \bm{X}_{b,0,k,t} = Y_r^{(b)}[k,t],
    \qquad
    \bm{X}_{b,1,k,t} = Y_i^{(b)}[k,t].
\end{equation}

Target-ul modelului este:
\begin{equation}
    \bm{M}^{\star}_{b,0,k,t} = M_r^{\star (b)}[k,t],
    \qquad
    \bm{M}^{\star}_{b,1,k,t} = M_i^{\star (b)}[k,t].
\end{equation}

\section{Integrarea VAD soft}

Pentru presetul LibriSpeech recomandat, proiectul foloseste etichete VAD soft:
\begin{equation}
    v[t] \in [0,1].
\end{equation}

Acestea sunt generate cu detectorul adaptiv VAD si salveaza scoruri continue de
incredere pentru prezenta vorbirii. Pentru a evita anularea completa a zonelor
ambigue, se foloseste un prag minim de trecere:
\begin{equation}
    g[t] = f + (1-f)v[t],
\end{equation}
unde $f$ este \textit{soft mask floor}. In presetul curent:
\begin{equation}
    f = 0.15.
\end{equation}

Pentru Complex MaskNet, gating-ul VAD se aplica pe ambele canale ale mastii CRM:
\begin{equation}
    M_r^{\star}[k,t] \leftarrow g[t]M_r^{\star}[k,t],
    \qquad
    M_i^{\star}[k,t] \leftarrow g[t]M_i^{\star}[k,t].
\end{equation}

Astfel, modelul invata mai putin agresiv in zonele de pauza sau incertitudine,
dar nu pierde complet informatia din zonele de tranzitie.

\section{Arhitectura Complex MaskNet}

Modelul foloseste aceeasi structura de baza ca MaskNet, dar cu intrare si iesire
complexe reprezentate prin doua canale reale.

\subsection{Configuratia generala}

Presetul principal este:
\begin{equation}
    \texttt{complex\_balanced\_res}.
\end{equation}

Configuratia sa este:
\begin{center}
\begin{tabular}{ll}
\toprule
Parametru & Valoare \\
\midrule
Canale intrare & 2 \\
Canale iesire & 2 \\
Base channels & 32 \\
Blocuri reziduale & Da \\
LSTM bidirectional & Da \\
Channel attention & Da \\
Activare iesire & $\tanh$ \\
Scalare iesire & 5.0 \\
Tip masca & Complex ratio mask \\
\bottomrule
\end{tabular}
\end{center}

\subsection{Bloc convolutional}

Blocul convolutional de baza este:
\begin{equation}
    h = \mathrm{ReLU}(\mathrm{BatchNorm}(\mathrm{Conv2D}(x))).
\end{equation}

\subsection{Bloc rezidual}

In varianta recomandata se folosesc blocuri reziduale:
\begin{align}
    z_1 &= F_1(x), \\
    z_2 &= F_2(z_1), \\
    h &= \mathrm{ReLU}(z_2 + P(x)),
\end{align}
unde $P(x)$ este fie identitatea, fie o proiectie $1 \times 1$ daca numarul de
canale se modifica.

Conexiunile reziduale ajuta propagarea gradientilor si stabilizeaza antrenarea
modelului mai adanc.

\subsection{Encoder}

Encoder-ul are patru niveluri. Fie $h_0 = X_{\mathrm{in}}$. Pentru fiecare nivel:
\begin{align}
    e_i &= \mathrm{ResidualBlock}_i(h_{i-1}), \\
    h_i &= \mathrm{MaxPool2D}(e_i).
\end{align}

Tensorii $e_i$ sunt pastrati pentru conexiunile de tip skip din decoder.

\subsection{Bottleneck si BiLSTM}

Dupa encoder, modelul aplica un bottleneck convolutional:
\begin{equation}
    z = \mathrm{Bottleneck}(h_4).
\end{equation}

Pentru modelarea dependintelor temporale, tensorul este rearanjat si introdus
intr-un LSTM bidirectional:
\begin{equation}
    z_{\mathrm{lstm}} = \mathrm{BiLSTM}(z).
\end{equation}

Rolul BiLSTM este sa modeleze contextul temporal al vorbirii, deoarece un cadru
STFT nu este independent de cadrele vecine.

\subsection{Channel attention}

Pe conexiunile skip se aplica atentie pe canale. Pentru un tensor $e$:
\begin{equation}
    a = \sigma(W_2 \delta(W_1 \mathrm{GAP}(e))),
\end{equation}
unde:
\begin{itemize}
    \item $\mathrm{GAP}$ este global average pooling;
    \item $\delta$ este activarea ReLU;
    \item $\sigma$ este sigmoid;
    \item $a$ sunt ponderile pe canale.
\end{itemize}

Feature-urile recalibrate sunt:
\begin{equation}
    e' = a \odot e.
\end{equation}

\subsection{Decoder}

Decoder-ul reconstruieste rezolutia timp-frecventa prin upsampling:
\begin{equation}
    u_i = \mathrm{ConvTranspose2D}(h_{i+1}).
\end{equation}

Apoi concateneaza skip connection-ul corespunzator:
\begin{equation}
    d_i = \mathrm{ResidualBlock}([u_i, e'_i]).
\end{equation}

\subsection{Stratul de iesire}

Ultimul strat produce doua canale:
\begin{equation}
    Z = \mathrm{Conv2D}_{1 \times 1}(d_1),
    \qquad
    Z \in \mathbb{R}^{B \times 2 \times K \times T}.
\end{equation}

Pentru ca CRM este limitata in intervalul $[-5,5]$, iesirea foloseste:
\begin{equation}
    \hat{M} = c \cdot \tanh(Z),
\end{equation}
unde $c=5.0$.

Astfel, modelul poate prezice valori negative sau pozitive pentru componentele
reala si imaginara ale mastii complexe, dar ramane stabil numeric.

\section{Functia de pierdere}

Pentru Complex MaskNet se foloseste MSE intre masca complexa prezisa si target-ul
CRM:
\begin{equation}
    \mathcal{L}_{\mathrm{CRM}} =
    \frac{1}{2BKT}
    \sum_{b=1}^{B}
    \sum_{c=1}^{2}
    \sum_{k=1}^{K}
    \sum_{t=1}^{T}
    \left(
        \hat{M}_{b,c,k,t} - M^{\star}_{b,c,k,t}
    \right)^2.
\end{equation}

Aceasta pierdere optimizeaza direct estimarea mastii complexe. In contrast,
MaskNet clasic optimizeaza o masca reala aplicata pe magnitudine.

\section{Pipeline de pregatire a datelor}

Pipeline-ul de date pentru presetul LibriSpeech Complex MaskNet este:

\begin{enumerate}
    \item Se incarca perechile clean/noisy.
    \item Audio este convertit la mono daca este necesar.
    \item Audio este resamplat la 16 kHz daca este necesar.
    \item Semnalele clean si noisy sunt aliniate la aceeasi lungime.
    \item Se aplica optional crop la lungimea maxima configurata.
    \item Se calculeaza STFT complex pentru clean:
    \begin{equation}
        S = \mathrm{STFT}(s).
    \end{equation}
    \item Se calculeaza STFT complex pentru noisy:
    \begin{equation}
        Y = \mathrm{STFT}(y).
    \end{equation}
    \item Se converteste $Y$ in doua canale: real si imaginar.
    \item Se calculeaza target-ul CRM $M^{\star}$.
    \item Daca VAD soft este activ, CRM este ponderata cu gate-ul VAD.
    \item Dataloader-ul returneaza:
    \begin{equation}
        \{\texttt{noisy\_complex}, \texttt{complex\_mask}\}.
    \end{equation}
\end{enumerate}

\section{Pipeline de antrenare}

Pentru fiecare batch:

\begin{enumerate}
    \item Se citeste intrarea:
    \begin{equation}
        \bm{X} = \texttt{batch["noisy\_complex"]}.
    \end{equation}
    \item Se citeste target-ul:
    \begin{equation}
        \bm{M}^{\star} = \texttt{batch["complex\_mask"]}.
    \end{equation}
    \item Modelul prezice masca:
    \begin{equation}
        \hat{\bm{M}} = f_{\theta}(\bm{X}).
    \end{equation}
    \item Se calculeaza pierderea:
    \begin{equation}
        \mathcal{L} =
        \mathrm{MSE}(\hat{\bm{M}}, \bm{M}^{\star}).
    \end{equation}
    \item Se aplica backpropagation.
    \item Se actualizeaza parametrii cu AdamW.
    \item Scheduler-ul OneCycleLR actualizeaza learning rate-ul.
    \item Se salveaza checkpoint-uri periodice si cel mai bun model pe validare.
\end{enumerate}

\section{Pipeline de inferenta}

La inferenta, pentru un semnal zgomotos $y[n]$:

\begin{enumerate}
    \item Se calculeaza spectrograma complexa:
    \begin{equation}
        Y = \mathrm{STFT}(y).
    \end{equation}
    \item Se creeaza tensorul cu doua canale:
    \begin{equation}
        \bm{X} =
        \begin{bmatrix}
            Y_r \\
            Y_i
        \end{bmatrix}.
    \end{equation}
    \item Modelul prezice CRM:
    \begin{equation}
        \hat{M} = f_{\theta}(\bm{X}).
    \end{equation}
    \item Se reconstruieste spectrograma complexa estimata:
    \begin{equation}
        \hat{S} = \hat{M} \cdot Y.
    \end{equation}
    \item Se aplica iSTFT:
    \begin{equation}
        \hat{s}[n] = \mathrm{iSTFT}(\hat{S}).
    \end{equation}
\end{enumerate}

Aceasta este diferenta esentiala fata de MaskNet clasic: output-ul final este
obtinut dintr-o spectrograma complexa estimata, nu dintr-o magnitudine estimata
combinata cu faza zgomotoasa initiala.

\section{Preseturi si comenzi de rulare}

\subsection{Generarea etichetelor VAD soft}

Pentru LibriSpeech + DEMAND, etichetele VAD soft se genereaza cu:

\begin{lstlisting}[language=bash]
python -m src.data_prep.generate_vad_labels --split all --label-set adaptive_soft_v1 --label-mode soft
\end{lstlisting}

Acestea sunt salvate in:

\begin{lstlisting}[language=bash]
data/librispeech_demand/vad_labels/adaptive_soft_v1/
\end{lstlisting}

\subsection{Antrenare Complex MaskNet pe LibriSpeech}

Presetul recomandat este:

\begin{lstlisting}[language=bash]
python -m src.train.train_mask_model --experiment-preset librispeech_complex_masknet_recommended
\end{lstlisting}

Pentru test rapid:

\begin{lstlisting}[language=bash]
python -m src.train.train_mask_model --experiment-preset librispeech_complex_masknet_recommended --max-batches 30
\end{lstlisting}

\subsection{Antrenare Complex MaskNet pe LibriSpeech}

Pentru faza de antrenare si validare se foloseste LibriSpeech-DEMAND:

\begin{lstlisting}[language=bash]
python -m src.train.train_mask_model --experiment-preset librispeech_complex_masknet_recommended
\end{lstlisting}

\subsection{Dry-run pentru verificarea configuratiei}

\begin{lstlisting}[language=bash]
python -m src.train.train_mask_model --experiment-preset librispeech_complex_masknet_recommended --dry-run
\end{lstlisting}

\section{Fisiere implementate in proiect}

\begin{center}
\begin{tabular}{ll}
\toprule
Fisier & Rol \\
\midrule
\texttt{src/dsp/stft\_utils.py} & Conversii complex-real/imag, CRM, iSTFT complex \\
\texttt{src/data\_prep/dataset.py} & Genereaza \texttt{noisy\_complex} si \texttt{complex\_mask} \\
\texttt{src/models/mask\_model.py} & MaskNet configurabil pentru output complex \\
\texttt{src/models/inference.py} & Inferenta diferentiata pentru masca reala sau complexa \\
\texttt{src/train/train\_loop.py} & Selecteaza target-ul potrivit: IRM sau CRM \\
\texttt{src/config.py} & Preseturi pentru Complex MaskNet \\
\bottomrule
\end{tabular}
\end{center}

\section{Comparatie intre cele doua variante}

\begin{center}
\begin{tabular}{lll}
\toprule
Caracteristica & MaskNet clasic & Complex MaskNet \\
\midrule
Intrare & $|Y|$ & $(Y_r, Y_i)$ \\
Canale intrare & 1 & 2 \\
Target & IRM & CRM \\
Canale iesire & 1 & 2 \\
Activare iesire & Sigmoid & Tanh scalat \\
Interval iesire & $[0,1]$ & $[-5,5]$ \\
Faza & Refoloseste faza noisy & Corecteaza implicit faza \\
Reconstructie & Magnitudine + faza noisy & Spectrograma complexa estimata \\
\bottomrule
\end{tabular}
\end{center}

\section{Avantaje si limitari}

\subsection{Avantaje}

\begin{itemize}
    \item Modelul poate modifica atat magnitudinea, cat si faza semnalului.
    \item Reconstructia audio este realizata dintr-o spectrograma complexa estimata.
    \item CRM ofera o formulare mai expresiva decat IRM.
    \item Arhitectura pastreaza avantajele MaskNet: U-Net, skip connections,
    BiLSTM si attention.
    \item Varianta clasica ramane disponibila ca baseline experimental.
\end{itemize}

\subsection{Limitari}

\begin{itemize}
    \item Target-ul CRM poate avea valori mari in zone cu energie scazuta, de aceea
    este necesar clipping.
    \item Modelul are intrare si iesire cu doua canale, deci foloseste mai multa
    memorie decat MaskNet clasic.
    \item Antrenarea poate fi mai sensibila la normalizare si la calitatea perechilor
    clean/noisy.
    \item Metricile obiective trebuie comparate cu baseline-ul clasic pentru a valida
    castigul real.
\end{itemize}

\section{Concluzie}

Complex MaskNet extinde sistemul de speech enhancement prin trecerea de la o masca
reala pe magnitudine la o masca complexa aplicata pe spectrograma STFT completa.
Prin estimarea componentelor reala si imaginara ale CRM, modelul poate influenta
atat amplitudinea, cat si faza semnalului reconstruit. Aceasta transformare face
pipeline-ul mai expresiv si mai potrivit pentru scenarii in care refolosirea fazei
zgomotoase limiteaza calitatea perceptuala a output-ului curat.

\end{document}
